\chapter{Implementace}
Tato kapitola navazuje na návrh řešení uvedený v předchozí kapitole a popisuje konkrétní realizaci navrženého systému. Zatímco kapitola~6 vymezila architekturu řešení, modelovou silniční síť, zónování, simulační metodiku a způsob vyhodnocení scénářů uzavírek, tato kapitola se zaměřuje na jejich převod do implementované výpočetní pipeline, datových struktur a uživatelské aplikace.

Implementace je založena na modulární pipeline řízené skriptem \texttt{run.py}. Jednotlivé kroky pipeline zajišťují přípravu silniční sítě, tvorbu dopravních okrsků, generování poptávky, dopravní přiřazení, kalibraci, validaci a publikaci výsledků. Důraz je kladen zejména na tok dat mezi jednotlivými částmi systému, reprodukovatelnost výpočtu a oddělení obecné metodiky od městsky specifické konfigurace. Kapitola postupuje od technického prostředí a struktury projektu přes výpočetní části modelu až po webovou aplikaci, která výsledky zpřístupňuje uživateli.


\section{Výpočetní pipeline modelu}
\label{sec:impl-pipeline}

Implementace je postavena na výpočetním backendu v~jazyce Python~3.12. Pro správu dopravního modelu, práci s~O--D maticemi a výpočet statického dopravního přiřazení je použita knihovna \texttt{AequilibraE}. Pro prostorové operace nad síťovými a zónovými daty jsou využity knihovny \texttt{GeoPandas} a \texttt{Shapely}, pomocné grafové operace zajišťuje \texttt{NetworkX}. Výsledky modelu jsou zpřístupněny přes REST API vytvořené ve frameworku \texttt{FastAPI} a vizualizovány ve webové aplikaci postavené na \texttt{React}, \texttt{TypeScript} a mapové knihovně \texttt{Leaflet}. Prostorové výpočty jsou prováděny v~souřadnicovém systému S-JTSK (EPSG:5514).

Navržený systém je koncipován jako statické makroskopické přiřazení dopravy. Cílem tedy není simulovat pohyb jednotlivých vozidel ani časový vývoj front, ale vytvořit reprodukovatelný postup pro porovnání referenčního stavu sítě se scénáři dopravních uzavírek.

Výpočetní proces je realizován jako modulární pipeline spouštěná skriptem \texttt{run.py}. Jednotlivé kroky na sebe datově navazují a ukládají mezivýsledky, což umožňuje opakovat pouze vybrané části výpočtu bez nutnosti znovu vytvářet celý model.

\begin{lstlisting}[language=bash, caption={Spuštění kroku \texttt{build-network} pro model Brna.}, label={lst:run-build-network}]
python run.py --config config/brno/sim.yaml build-network
\end{lstlisting}

Z~hlediska funkce lze pipeline rozdělit do tří hlavních částí. První část připravuje silniční síť a zónování, druhá část vytváří dopravní poptávku a O--D matici a třetí část provádí dopravní přiřazení, kalibraci, validaci a zpřístupnění výsledků pro scénářové výpočty. Hlavním integračním bodem je projektová databáze knihovny \texttt{AequilibraE}, do které se ukládá síťový model, matice i výstupy přiřazení.

\begin{table}[ht]
\centering
\small
\begin{tabular}{rll}
\toprule
\textbf{\#} & \textbf{Krok} & \textbf{Popis} \\
\midrule
0  & \texttt{init-city}         & Inicializace konfigurace a~adresářů pro nové město \\
1  & \texttt{clean}             & Smazání vygenerovaných výstupů \\
2  & \texttt{check}             & Ověření integrity AequilibraE projektu \\
3  & \texttt{build-network}     & Import silniční sítě z~OpenStreetMap \\
4  & \texttt{normalize-network} & Normalizace atributů sítě a~aplikace baseline uzavírek \\
5  & \texttt{build-zones}       & Tvorba dopravních zón, centroidů a~konektorů \\
6  & \texttt{fetch-data}        & Stažení externích datových sad (SLDB, CSD, RÚIAN, uzavírky) \\
7  & \texttt{build-supernetwork}& Hrubý národní graf pro klasifikaci externích toků \\
8  & \texttt{build-demand}      & Sestavení výchozí O-D matice ze všech segmentů poptávky \\
9  & \texttt{assign-warm-skims} & Předběžné přiřazení pro odhad impedanční matice \\
10 & \texttt{distribute}        & Gravitační distribuce a~IPF vyvažování matice \\
11 & \texttt{assign}            & Hlavní dopravní přiřazení algoritmem BFW \\
12 & \texttt{audit-supply}      & Diagnostika nabídkových parametrů sítě \\
13 & \texttt{calibrate}         & Vícefázová kalibrace modelu vůči sčítacím datům \\
14 & \texttt{calibrate-odme}    & ODME kalibrace (alternativa ke~\texttt{calibrate}) \\
15 & \texttt{tune-supply}       & Doladění nabídkových parametrů sítě \\
16 & \texttt{validate}          & Nezávislá validace vůči datům nepoužitým při kalibraci \\
17 & \texttt{sensitivity}       & Citlivostní analýza modelu \\
18 & \texttt{learn-profile}     & Učení časových profilů z~dat sčítání dopravy \\
19 & \texttt{strip-closures}    & Obnovení čistého stavu sítě před scénáři \\
20 & \texttt{serve}             & Spuštění REST API serveru s~výsledky modelu \\
\bottomrule
\end{tabular}
\caption{Kroky výpočetní pipeline spouštěné skriptem \texttt{run.py}.}
\label{tab:pipeline-steps}
\end{table}

Konfigurace modelu je oddělena od zdrojového kódu. Obecné parametry, jako jsou výchozí rychlosti, kapacity, parametry BPR funkcí nebo nastavení přiřazení, jsou definovány globálně ve frameworku. Městská konfigurace obsahuje pouze lokální údaje -- název modelované oblasti, zdroje zónování, filtr dojížďkových dat a přiřazení kalibračních dat. Při spuštění výpočtu je městská konfigurace automaticky sloučena s globálními výchozími hodnotami. Díky tomu lze stejnou pipeline použít i pro jiné město bez úprav implementace -- stačí vytvořit nový konfigurační adresář \texttt{config/<město>/} s minimálním souborem \texttt{sim.yaml}.

\begin{lstlisting}[language={},caption={Minimální městská konfigurace modelu (soubor \texttt{sim.yaml}).},label={lst:city-config}]
osm:
  place_name: "Olomouc, Czechia"

zoning:
  sources:
    - type: osm
      place: "Okres Olomouc, Czechia"
      admin_level: "9"

calibration:
  count_source: "csd_split"

datasets:
  sources:
    commuting_sldb2021:
      filter:
        origin:
          keep_if_any_matches:
            - field: "op_okres"
              values: ["Olomouc"]
\end{lstlisting}

Při zachování stejných vstupních dat a konfigurace systém generuje stejnou síťovou topologii, což je důležité pro porovnání referenčního stavu se scénáři dopravních uzavírek.


% =========================================================================
\section{Konstrukce referenčního modelu}
\label{sec:impl-network}

Tato sekce popisuje konstrukci referenčního modelu, tedy silniční sítě, dopravních zón a externích bran. Výsledný model města Brna sestává z~\textbf{52\,493~hran}, \textbf{47\,181~uzlů} a~\textbf{219~dopravních zón} (191~interních a~28~externích gateway zón).

\subsection{Import sítě z OpenStreetMap}
\label{sec:impl-network-osm}

Proces importu je zajištěn funkcí \texttt{build\_network\_from\_osm} v modulu \texttt{network/pipeline.py}. Vstupními parametry jsou identifikátor lokality (\texttt{osm.place\_name}) a volitelná šířka obalové zóny (\texttt{osm.buffer\_km}). V případě definování nenulového bufferu je nejprve provedeno geokódování názvu lokality pomocí knihovny \texttt{OSMnx}, výsledný polygon je metrickou projekcí rozšířen o zadanou vzdálenost a doplněn o 5\,km marži pro zajištění konektivity při stahování dat. Následně dochází k ořezu sítě zpět na hranice modelového polygonu, čímž je zamezeno nadbytečnému šíření dálničních koridorů mimo zájmovou oblast.

Vlastní pořízení dat je realizováno voláním metody \texttt{create\_from\_osm} rozhraní \texttt{AequilibraE}, která zprostředkovává získání orientovaného grafu ze serveru Overpass API. Bezprostředně po importu jsou aplikovány čtyři čisticí operace:

\begin{enumerate}
  \item \textbf{Filtrace nesjízdných hran} -- z grafu jsou odstraněny cesty určené výhradně pro pěší, cyklisty či specifické druhy dopravy (např.\ \texttt{footway}, \texttt{cycleway}, \texttt{pedestrian}). Celkem je filtrováno 13 typů komunikací. Zároveň jsou odstraněny hrany, které v atributu \texttt{modes} neobsahují příznak pro osobní automobily (\texttt{c}).
  \item \textbf{Extrakce největší souvislé komponenty} -- z neorientovaného grafu je pomocí knihovny \texttt{NetworkX} sestaven multigraf, ze kterého je zachována pouze největší komponenta. Tímto je zajištěna matematická řešitelnost přiřazení odstraněním izolovaných subgrafů.
  \item \textbf{Urban-core trim} -- vypočítá se bounding box městského jádra z~uzlů připojených k~hranám, které nejsou koridorového typu (\texttt{motorway}, \texttt{trunk}, \texttt{primary} a~jejich \texttt{*\_link} varianty). Nekoridorové hrany ležící mimo tento obdélník jsou odstraněny, zatímco páteřní koridory jsou zachovány vždy, protože na jejich koncích budou umístěny gateway uzly. Po~ořezu se znovu spustí filtr souvislosti.
  \item \textbf{OSM enrichment} -- pro každou hranu sítě jsou z nezjednodušeného grafu \texttt{OSMnx} zpětně doplněny atributy \texttt{osm\_ref} (číslo silnice) a \texttt{osm\_name}. Implementován je rovněž algoritmus pro propagaci referenčních čísel na jmenovaných koridorech, který vyplňuje krátké mezery v datech OSM na základě shody v okolních hranách.
\end{enumerate}

Geometrická data jsou vedena v souřadnicovém systému S-JTSK (EPSG:5514). Rozložení hran výsledného modelu Brna dle typu komunikace je uvedeno v tabulce~\ref{tab:link-types}.

\begin{table}[ht]
\centering
\label{tab:link-types}
\begin{tabular}{lr}
\toprule
\textbf{Typ komunikace} & \textbf{Počet hran} \\
\midrule
residential    & 21\,101 \\
service        & 15\,851 \\
tertiary       & 6\,365  \\
secondary      & 3\,792  \\
living\_street & 1\,905  \\
unclassified   & 1\,005  \\
trunk          & 654     \\
primary        & 533     \\
trunk\_link    & 454     \\
motorway       & 356     \\
motorway\_link & 223     \\
ostatní        & 254     \\
\midrule
\textbf{Celkem} & \textbf{52\,493} \\
\bottomrule
\end{tabular}
\caption{Počet hran modelu Brna dle typu komunikace.}
\end{table}

\subsection{Normalizace atributů}
\label{sec:impl-network-norm}

Vzhledem k neúplnosti surových dat OSM v oblastech rychlostních limitů či počtu jízdních pruhů je aplikován krok \texttt{normalize-network}. Směrové atributy (\texttt{*\_ab}, \texttt{*\_ba}) jsou zpracovávány odděleně pro zachování směrové asymetrie. Algoritmus využívá hierarchické doplňování ve třech vrstvách:

\begin{enumerate}
  \item \textbf{Data OSM} -- prioritně jsou využity hodnoty přímo z tagů (např.\ \texttt{maxspeed}).
  \item \textbf{Kopie protisměru} -- u obousměrných hran je chybějící hodnota v jednom směru nahrazena známou hodnotou ze směru opačného.
  \item \textbf{Výchozí parametry frameworku} -- zbývající hodnoty jsou doplněny ze slovníku \texttt{NETWORK\_NORM\_DEFAULTS} na základě typu komunikace. Odhadnuté hodnoty jsou indikovány příznakem \texttt{estimated\_*}, což umožňuje jejich odlišení při diagnostice modelu.
\end{enumerate}

Kromě doplňování chybějících hodnot normalizace aplikuje rovněž konfigurovatelné experimentální profily (\texttt{network\_normalization.yaml}), které umožňují překlasifikovat typy hran, nastavit stropy a~podlahy rychlostí nebo upravit kapacitní faktory pro konkrétní koridory. Po~normalizaci atributů se provádějí opravy konektivity sítě: jednosměrné dálniční hrany na okraji modelu bez protisměrného pásu jsou převedeny na obousměrné a~opravují se slepé konce dělených komunikací.

Délka hrany se počítá z její geometrie v metrech. Na základě délky, volné rychlosti a počtu pruhů jsou následně odvozeny základní atributy potřebné pro přiřazení dopravy. Výpočet volného cestovního času, směrové kapacity a minimálních prahů je shrnut ve vztahu~\ref{eq:network-basic-attributes}.

\begin{equation}
\begin{aligned}
t_0 &= \frac{d}{v_0 / 3{,}6}, \\
c &= n_{\mathrm{lanes}} \cdot c_{\mathrm{lane}}, \\
v_0 &\geq 5~\mathrm{km/h}, \qquad
c \geq 50~\mathrm{voz/hod}, \qquad
t_0 \geq 0{,}01~\mathrm{s}.
\end{aligned}
\label{eq:network-basic-attributes}
\end{equation}

Ve vztahu značí $t_0$ volný cestovní čas (\emph{free-flow time}), $d$ délku hrany v metrech, $v_0$ volnou rychlost v km/h, $n_{\mathrm{lanes}}$ počet pruhů a $c_{\mathrm{lane}}$ kapacitu jednoho pruhu. Převod rychlosti výrazem $v_0 / 3{,}6$ převádí rychlost z km/h na m/s, takže výsledný čas vychází v sekundách.

\subsection{Parametrizace BPR funkcí}
\label{sec:impl-network-bpr}

Odporová funkce BPR (Bureau of Public Roads) je parametrizována specificky pro jednotlivé typy komunikací. Hodnoty koeficientů $\alpha$ a $\beta$ jsou uloženy v globální konfiguraci frameworku (viz tabulka~\ref{tab:bpr-params}). 

Zvolená metodika předpokládá nižší hodnotu $\alpha$ u dálničních koridorů (0,15), což reflektuje nižší citlivost víceproudých komunikací na nárůst intenzit. Vyšší hodnoty u lokálních komunikací (až 1,00) pak modelují strmější nárůst zpoždění při překročení kapacity. Denní kapacitní faktor $f_\mathrm{daily}$ slouží k přepočtu hodinové kapacity na denní, přičemž vyšší hodnoty u páteřních komunikací zohledňují jejich rovnoměrnější denní zatížení.

\begin{table}[ht]
\centering
\label{tab:bpr-params}
\begin{tabular}{lrrr}
\toprule
\textbf{Typ komunikace} & $\alpha$ & $\beta$ & \textbf{Denní kap. faktor} \\
\midrule
motorway       & 0,15 & 4,0 & 13,0 \\
motorway\_link & 0,25 & 4,0 & 11,0 \\
trunk          & 0,45 & 4,0 & 12,0 \\
trunk\_link    & 0,50 & 4,0 & 10,0 \\
primary        & 0,60 & 4,0 & 10,0 \\
secondary      & 0,70 & 4,0 & 9,0  \\
tertiary       & 0,75 & 4,0 & 8,0  \\
residential    & 0,85 & 4,0 & 6,5  \\
living\_street & 1,00 & 4,0 & 5,5  \\
service        & 1,00 & 4,0 & 5,5  \\
\bottomrule
\end{tabular}
\caption{BPR parametry $\alpha$, $\beta$ a~denní kapacitní faktor (framework defaults).}
\end{table}

Tyto hodnoty byly zvoleny jako výchozí heuristické nastavení a slouží především jako počáteční parametrizace před kalibrací. Nejde tedy o samostatně validované hodnoty pro město Brno.

\subsection{Zapracování historických uzavírek}
\label{sec:impl-network-closures}

Za účelem dosažení shody baseline přiřazení s reálnou situací v období kalibrace jsou do sítě promítnuty historické uzavírky. Vstupem je datová sada ve formátu Parquet obsahující prostorové a časové atributy omezení. Uzavírky jsou filtrovány dle období definovaného parametrem \texttt{measurement\_period}.

Párování uzavírek na hrany sítě je realizováno prostorovou analýzou (vytvořením bufferu o poloměru 150\,m kolem bodu uzavírky) s volitelnou validací čísla silnice (\texttt{road\_ref}). Na dotčené hrany jsou následně aplikovány redukční faktory pro kapacitu a rychlost v závislosti na závažnosti omezení (\textit{full}, \textit{lane reduction}, \textit{speed limit}). Původní hodnoty jsou archivovány v pomocných atributech pro možnost následné obnovy sítě do čistého stavu v kroku \texttt{strip-closures}.

\subsection{Exportní formáty}
\label{sec:impl-network-export}

V závěru normalizace je síť exportována do tří formátů pro různé účely využití:
\begin{itemize}
    \item \textbf{GeoPackage (GPKG)} -- v metrickém souřadnicovém systému S-JTSK pro vnitřní výpočty modelu.
    \item \textbf{GeoJSON} -- v souřadnicovém systému WGS84 pro účely vizualizace ve webovém rozhraní.
    \item \textbf{Parquet} -- pro efektivní analytické dotazy bez prostorové složky.
\end{itemize}
Všem entitám jsou přiřazeny stabilní identifikátory (\texttt{link\_id}, \texttt{node\_id}), které zajišťují konzistenci dat napříč všemi moduly systému, od výpočetního jádra až po frontendovou vizualizaci.


% =========================================================================
\section{Zónování a napojení poptávky na síť}
\label{sec:impl-zones}

Proces zónování, realizovaný modulem \texttt{build-zones}, zahrnuje rozdělení modelového území na dopravní okrsky (Traffic Analysis Zones, TAZ), jejich napojení na fyzickou silniční síť pomocí centroidových konektorů a identifikaci vstupních bodů pro vnější dopravní vztahy. Model města Brna v rámci této práce využívá celkem \textbf{219 zón}, z nichž 191 tvoří interní okrsky a 28 reprezentuje externí gateway zóny pro modelování vjezdu a výjezdu z oblasti.

\subsection{Definice dopravních okrsků}
\label{sec:impl-zones-definition}

Dopravní zóny jsou definovány prostřednictvím konfigurovatelného systému zdrojů v modulu \texttt{zoning/pipeline.py}. Navržený framework podporuje dva základní režimy definice územního členění:
\begin{enumerate}[label=(\alph*)]
    \item automatické stažení administrativních hranic z databáze OpenStreetMap na základě definované úrovně (\texttt{admin\_level}),
    \item import externích prostorových dat ve formátu GeoJSON s vlastním členěním.
\end{enumerate}

Jednotlivé zdroje lze kombinovat v hierarchické struktuře, kde parametr \texttt{source\_rank} určuje prioritu při automatickém odstraňování překryvů ploch. V případě modelu Brna je primárním zdrojem (\texttt{source\_rank=0}) 76 katastrálních území města, která jsou doplněna o 115 okolních obcí ze sekundárního zdroje (\texttt{source\_rank=1}).

Algoritmus zpracování postupuje v následujících krocích:
\begin{itemize}
    \item \textbf{Načtení a filtrace:} Jsou načteny geometrie polygonů, přičemž jsou uvažovány pouze zóny, jejichž reprezentativní bod leží uvnitř zájmové oblasti (Area of Interest, AOI). AOI je definována jako konvexní obálka silniční sítě s aplikovanou kvantilovou ořezkou (1\,\%) pro eliminaci odlehlých uzlů.
    \item \textbf{Topologické vyčištění:} Dochází k odečtení překryvů ploch na základě priorit zdrojů, čímž je zajištěno, že každý bod prostoru náleží právě jedné zóně.
\end{itemize}

Výsledkem je vrstva \texttt{zones.geojson}, která nese nezbytná metadata pro navazující výpočty poptávky a identifikaci tranzitních uzlů.

\subsection{Centroidy a konektory}
\label{sec:impl-zones-centroids}

Aby bylo možné přiřadit dopravní poptávku na fyzickou síť, je pro každou zónu vygenerován právě jeden centroidový uzel. Tento uzel je v databázi \texttt{AequilibraE} registrován s příznakem \texttt{is\_centroid=1}. U interních zón je poloha centroidu ztotožněna s reprezentativním bodem polygonu, zatímco u externích zón je centroid umístěn vně hranic modelu, aby konektory korektně reprezentovaly příjezdové trasy z vnějšího prostředí.

Napojení na silniční síť zajišťuje algoritmus \texttt{create\_centroid\_connectors}, který postupuje podle následujících pravidel:
\begin{enumerate}
    \item \textbf{Výběr kandidátních uzlů:} Jsou uvažovány pouze uzly náležející do největší silně souvislé komponenty (SCC) s povoleným automobilovým režimem.
    \item \textbf{Hierarchická filtrace:} Z výběru jsou exkludovány uzly připojené výhradně k dálniční síti (\texttt{motorway}, \texttt{motorway\_link}). Tímto opatřením je zamezeno nereálnému generování dopravy přímo na dálniční koridory bez využití sjezdů a komunikací nižších tříd.
    \item \textbf{Vytvoření konektorů:} Pro každou zónu je vytvořeno až šest konektorů k nejbližším vhodným uzlům v dosahu 3\,000\,m. Výběr je vážen třídou komunikace (\texttt{road\_weight}) a směrovou diverzitou.
    \item \textbf{Parametrizace:} Interní konektory jsou definovány jako obousměrné hrany s~rychlostí 20\,km/h a~časovou penalizací 300\,s. Tato penalizace (přístupový odpor) zabraňuje tomu, aby algoritmus přiřazení využíval centroidy jako zkratky. Konektory externích gateway zón mají odlišné parametry: rychlost 90\,km/h a~penalizaci pouhých 5\,s, aby realisticky reprezentovaly vjezd z~páteřních komunikací.
\end{enumerate}

\subsection{Supernetwork a externí brány}
\label{sec:impl-zones-supernetwork}

Pro realistickou reprezentaci vnějších dopravních vztahů byl implementován modul \texttt{supernetwork.py}. Jeho úkolem je identifikovat vstupní a výstupní brány (\textit{gateways}) na hranici modelové oblasti a klasifikovat relace mezi externími obcemi. Výpočet je rozdělen do čtyř navazujících fází:

\begin{enumerate}
    \item \textbf{Konstrukce národního grafu.} Z celostátních dat OpenStreetMap je extrahován zjednodušený graf páteřních komunikací. Pomocí grafové kontrakce jsou odstraněny nadbytečné uzly, přičemž zůstává zachována topologie hlavních dopravních koridorů.

    \item \textbf{Integrace socioekonomických dat.} Data o dojížďce ze sčítání lidu (SLDB 2021) jsou napárována na uzly národního grafu. V případě modelu Brna bylo identifikováno 6\,117 externích obcí s kompletním pokrytím dojížďkových proudů.

    \item \textbf{Přiřazení bran.} Pro každou vnější relaci je pomocí Dijkstrova algoritmu nalezeno vhodné spojení přes hranici modelové oblasti. V modelu Brna bylo identifikováno 9 klíčových gateway bodů, mezi které patří například dálniční tahy D1 ve směru na Prahu, D2 ve směru na Bratislavu a silnice I/43 ve směru na Svitavy.

    \item \textbf{Klasifikace a filtrace tranzitu.} Relace jsou rozděleny podle toho, zda mají být v modelu reprezentovány jako externí nebo tranzitní. Tranzitní relace jsou zahrnuty pouze tehdy, pokud průjezd modelovaným územím představuje časově efektivní alternativu k jeho objetí. Pro tuto filtraci jsou použity dva současně platné parametry: \texttt{detour\_ratio\_max}~$= 1{,}25$ (poměr objížďky nesmí překročit 125\,\% přímé cesty) a~\texttt{max\_extra\_minutes}~$= 18$ (absolutní přírůstek nesmí překročit 18~minut). V modelu Brna bylo tímto způsobem akceptováno 26\,714 tranzitních párů, které generují denní objem přibližně 146\,982 vozidel.
\end{enumerate}

Výstupy této fáze jsou uloženy v lookup tabulkách ve formátu Parquet. Tyto tabulky následně slouží jako vstup pro konstrukci matic dopravní poptávky v navazujících krocích pipeline.


\section{Generování dopravní poptávky}
\label{sec:impl-demand}

Generování dopravní poptávky je realizováno jako samostatný krok pipeline (\texttt{build-demand}), který na základě heterogenních datových zdrojů sestaví kompletní denní O-D matici pro modelované území. Výsledkem je strukturovaná matice rozdělená na jednotlivé segmenty (commuting, other, external) a časové periody, která slouží jako vstup pro přiřazovací algoritmus.

\subsection{Datové zdroje}
\label{sec:impl-demand-sources}

Vstupní data zajišťuje krok \texttt{fetch-data}, který stahuje a~konvertuje datové sady do~jednotného formátu Parquet. Hlavním zdrojem jsou meziobecní dojížďkové relace ze SLDB~2021 (filtrované na lokalizace \texttt{0\_na\_adrese\_OP} a~\texttt{1\_meziobecni}), kalibrační sčítací profily CSD (69~profilů, z~nichž 44 je po~filtraci použito) a~centroidy obcí z~registru RÚIAN pro geografické ukotvení zón.

\subsection{Dojížďkové jádro (commuting core)}
\label{sec:impl-demand-commuting}

Hlavní část O-D matice je odvozena z dojížďkových dat SLDB. Konstrukce tohoto segmentu je realizována funkcí \texttt{load\_or\_build\_od\_matrix} v modulu \texttt{demand/pipeline.py}. V případě modelu Brna je vstupem dataset obsahující 253\,377 relací mezi obcemi.

Prvním krokem je mapování názvů obcí na interní identifikátory dopravních zón (\texttt{zone\_id}). Tento proces využívá víceúrovňové párování, které kombinuje přímé shody, skupinové shody přes huby a externí mapování pomocí tabulky \texttt{gateway\_lookup.parquet}. Pro model Brna bylo tímto postupem nalezeno 41\,496 přímých shod, 4\,551 skupinových shod a 35\,573 externích mapování. Výsledkem je množina 40\,810 unikátních O-D párů použitých při konstrukci dojížďkového segmentu.

Po přiřazení relací k zónám jsou počty osob převedeny na vozidla pomocí konverzních parametrů definovaných v konfiguraci frameworku. Pro pracovní dojížďku jsou v modelu Brna použity hodnoty \texttt{car\_share = 0{,}48}, \texttt{occupancy = 1{,}20} a \texttt{trips\_per\_person = 2{,}0}. Pro školní dojížďku jsou použity hodnoty \texttt{car\_share = 0{,}25} a \texttt{occupancy = 1{,}30}. Tyto koeficienty převádějí počet dojíždějících osob na odhadovaný počet jízd osobními vozidly.

Externí relace, tedy relace s počátkem nebo cílem mimo modelované území, jsou mapovány na gateway zóny pomocí lookup tabulky vytvořené v kroku supernetwork. Výsledný objem dojížďkového segmentu pro model Brna činí \textbf{691\,685 vozidel za den}.

\subsection{Gravitační model pro segment \uv{other}}
\label{sec:impl-demand-gravity}

Segment nepracovních cest (\uv{other}) je generován synteticky pomocí gravitačního modelu (viz kapitola~4). Implementace využívá exponenciální deterrenční funkci s parametrem $\beta = 0{,}00030$, který je definován ve frameworku (\texttt{demand.segments.other.beta}).

Produkce cest je odvozena z populace jednotlivých zón, přičemž výsledné objemy jsou převedeny na vozidla pomocí parametrů \texttt{car\_share = 0{,}35} a \texttt{occupancy = 1{,}50}. Model je aplikován pouze na interní zóny, externí gateway zóny jsou z tohoto segmentu vyloučeny.

Výsledný objem segmentu \uv{other} činí \textbf{135\,530 vozidel za den}.

\subsection{Externí toky}
\label{sec:impl-demand-external}

Externí doprava je do modelu doplněna ve formě dvou samostatných sub-segmentů, které vycházejí z výstupů supernetwork kroku.

První sub-segment (\textit{external-local}) reprezentuje vztahy mezi gateway zónami a interními zónami. Jeho objem je odvozen z dat sčítání dopravy a následně škálován parametrem \texttt{external\_commuting\_scale = 0{,}65}. V modelu Brna činí tento segment \textbf{35\,000 vozidel za den}.

Druhý sub-segment (\textit{external-through}) reprezentuje tranzitní dopravu mezi gateway zónami. Vstupem je matice \texttt{through\_gateway\_pairs.parquet}, která je výsledkem supernetwork analýzy. Tento segment je škálován faktorem \texttt{through\_traffic\_scale = 0{,}50} a jeho objem činí \textbf{73\,491 vozidel za den}.

Celkový objem externí dopravy je tedy \textbf{108\,491 vozidel za den}. Spolu s ostatními segmenty tvoří kompletní denní matici o objemu \textbf{935\,706 vozidel za den}.

\subsection{Časová disagregace}
\label{sec:impl-demand-temporal}

Denní O-D matice je rozdělena do čtyř časových period (AM, IP, PM, EV) podle parametrů definovaných ve frameworku (\texttt{demand.time\_slices}). Rozdělení probíhá samostatně pro každý segment, přičemž jsou zohledněny rozdílné časové profily jednotlivých typů cest.

Pracovní cesty jsou rozděleny asymetricky (outbound do ranní špičky, return do odpolední), zatímco nepracovní a tranzitní cesty mají rovnoměrnější distribuci.

Výsledná matice je uložena v nativním formátu AequilibraE (\texttt{.aem}) jako multi-class struktura se dvěma třídami:
\begin{itemize}
    \item \texttt{wd\_daily\_local} (lokální doprava): 862\,214 voz./den,
    \item \texttt{wd\_daily\_external\_through} (tranzitní doprava): 73\,491 voz./den.
\end{itemize}

Toto rozdělení umožňuje samostatné zpracování tranzitní a lokální dopravy při následném přiřazení (viz sekce~\ref{sec:impl-calibration-assignment}).


% =========================================================================
\section{Distribuce a kalibrace modelu}
\label{sec:impl-calibration}

Po sestavení počáteční O-D matice je nutné upravit poptávku tak, aby lépe odpovídala prostorové struktuře modelované sítě a dostupným pozorovaným intenzitám dopravy. Tato část pipeline proto zahrnuje distribuci poptávky, dopravní přiřazení a následnou kalibraci modelu vůči sčítacím profilům. Implementačně jsou tyto kroky realizovány zejména v modulech \texttt{distribution/pipeline.py}, \texttt{assignment/} a~\texttt{calibration/}.

Distribuce zpřesňuje mezizonální vztahy na základě impedance mezi zónami, dopravní přiřazení převádí O-D matici na zatížení jednotlivých hran sítě a kalibrace následně upravuje poptávku podle rozdílů mezi modelovanými a pozorovanými objemy. V této kapitole je popsán implementační mechanismus uvedených kroků. Konkrétní dosažené kalibrační a validační výsledky jsou vyhodnoceny až v kapitole~8.

\subsection{Gravitační distribuce a IPF vyvažování}
\label{sec:impl-calibration-distribution}

Počáteční matice vytvořená v kroku \texttt{build-demand} vychází z dojížďkových dat, syntetických segmentů a externích toků. U takto sestavené matice však nemusí být prostorové rozložení cest plně v souladu s reálnou dostupností mezi jednotlivými zónami. Krok \texttt{distribute} proto upravuje mezizonální vztahy pomocí gravitačního modelu a následného IPF vyvažování.

Vstupem distribuce je impedanční matice. Pokud byl předem spuštěn krok \texttt{assign-warm-skims}, použijí se jako impedance síťové cestovní časy uložené ve skim matici. Pokud tyto skimy dostupné nejsou, framework použije euklidovské vzdálenosti mezi centroidy zón. Tento záložní režim umožňuje spustit distribuci i v rané fázi sestavení modelu, avšak pro finální výpočet jsou vhodnější síťové skimy, protože lépe reprezentují skutečnou dostupnost v uliční síti.

Na základě impedanční matice je aplikována exponenciální deterrenční funkce gravitačního modelu popsaná v kapitole~4. Výsledná gravitační matice je následně vyvážena metodou IPF/Furness tak, aby odpovídala zadaným produkcím a atrakcím jednotlivých zón. V implementaci však tato matice zcela nenahrazuje původní seed matici. Obě matice jsou kombinovány váženým průměrem podle vztahu~\ref{eq:distribution-blend}.

\begin{equation}
T_{\mathrm{final}} = \alpha \cdot T_{\mathrm{IPF}} + (1 - \alpha) \cdot T_{\mathrm{seed}}
\label{eq:distribution-blend}
\end{equation}

Ve vztahu~\ref{eq:distribution-blend} označuje $T_{\mathrm{final}}$ výslednou distribuovanou matici, $T_{\mathrm{IPF}}$ gravitační matici po IPF vyvážení a $T_{\mathrm{seed}}$ původní počáteční matici z kroku \texttt{build-demand}. Parametr $\alpha$ určuje váhu IPF matice. Framework poskytuje výchozí hodnotu $\alpha = 0{,}7$, zatímco model Brna používá městský přepis $\alpha = 0{,}5$. Tento postup zachovává část struktury původních dojížďkových dat a současně zohledňuje síťovou impedanci mezi zónami.

\subsection{Dopravní přiřazení}
\label{sec:impl-calibration-assignment}

Dopravní přiřazení převádí O-D matici na zatížení jednotlivých hran modelové sítě. V implementaci je tento krok realizován modulem \texttt{assignment.py}. Pro hlavní výpočet je použit algoritmus BFW (Bi-conjugate Frank-Wolfe), který slouží k řešení rovnovážného přiřazení poptávky na síť popsaného v kapitole~4.

Náklad hrany je odvozen od cestovního času vypočteného pomocí BPR funkce a doplněn o distanční složku. Tato složka je řízena parametrem \texttt{fixed\_cost\_multiplier} a omezuje výběr nepřiměřeně dlouhých objížděk v případech, kdy by byly časově podobné kratší trase. Výstupem přiřazení je tabulka zatížení hran uložená do souboru \texttt{assignment\_results.parquet}. Volitelně mohou být uloženy také skim matice využitelné v navazujících krocích.

Přiřazení pracuje se dvěma třídami vozidel. Lokální a tranzitní doprava jsou přiřazovány současně, ale s odlišnými parametry generalizovaného nákladu. Nastavení tříd shrnuje tabulka~\ref{tab:multi-class}.

\begin{table}[ht]
\centering
\label{tab:multi-class}
\begin{tabular}{llrr}
\toprule
\textbf{Třída} & \textbf{O-D core} & \textbf{VOT} & \textbf{PCE} \\
\midrule
local   & wd\_daily\_local            & 1,0 & 1,0 \\
through & wd\_daily\_external\_through & 2,0 & 1,5 \\
\bottomrule
\end{tabular}
\caption{Třídy vozidel pro multi-class přiřazení.}
\end{table}

Vyšší hodnota VOT u tranzitní dopravy vyjadřuje větší citlivost na cestovní čas a tím i preferenci rychlejších kapacitních komunikací. Vyšší hodnota PCE u tranzitní třídy zohledňuje vyšší zastoupení těžších vozidel v tranzitním provozu. Oddělení tříd umožňuje pracovat s lokální a tranzitní dopravou odlišně, aniž by bylo nutné vytvářet samostatné síťové modely.

\subsection{Kalibrace O-D matice}
\label{sec:impl-calibration-odme}

Kalibrace slouží k úpravě O-D matice tak, aby modelované intenzity na vybraných sčítacích profilech lépe odpovídaly pozorovaným hodnotám. Kalibrační logika je realizována v modulech \texttt{calibration/odme.py} a~\texttt{calibration/legacy.py} a využívá výsledky dopravního přiřazení nad sítí AequilibraE. Základní kalibrační smyčka je založena na opakovaném přiřazení aktuální matice, porovnání modelovaných objemů s pozorovanými hodnotami a následné úpravě vybraných částí O-D matice.

Framework podporuje více typů korekcí. Globální škálování upravuje celou matici jedním koeficientem podle poměru celkových pozorovaných a modelovaných objemů, avšak nemění prostorový vzor poptávky. Entropijní aktualizace upravuje jednotlivé O-D páry multiplikativně a zachovává strukturu počáteční matice. Pro hlavní kalibraci je využita gradientová ODME varianta, která pracuje s informací o tom, které O-D páry přispívají k tokům na konkrétních sčítacích profilech.

Pro každý sčítací profil jsou pomocí select-link analýzy určeny proporce O-D párů, které přes daný profil projíždějí. Rozdíl mezi pozorovaným a modelovaným objemem je následně rozdělen zpět do těchto O-D párů podle jejich podílu na toku. Aby nedocházelo k nestabilním změnám, jsou korekce tlumeny damping parametry a výsledná matice je omezena vůči původní počáteční matici.

\subsection{Vícefázová kalibrační pipeline}
\label{sec:impl-calibration-multistage}

ODME kalibrace je citlivá na kvalitu počáteční matice, rozložení externích toků a prostorové pokrytí sčítacích profilů. Z tohoto důvodu je kalibrace implementována jako vícefázový proces, ve kterém jednotlivé fáze postupně řeší různé typy nesouladu mezi modelem a pozorovanými daty.

První fáze upravuje gravitační distribuci poptávky. Pokud jsou k dispozici síťové skimy, je znovu odhadnuta deterrenční funkce a matice je vyvážena metodou IPF. Cílem této fáze není měnit celkový objem poptávky, ale zpřesnit prostorové rozložení mezizonálních vztahů před samotnou kalibrací.

Druhá fáze se zaměřuje na gateway zóny. U profilů na vstupu do modelové oblasti se porovnají modelované a pozorované objemy a podle zjištěného poměru se upraví řádky a sloupce O-D matice vázané na příslušné gateway zóny. Tato předkalibrace pomáhá snížit systematické odchylky externí dopravy před detailnější ODME korekcí.

Třetí fázi tvoří vlastní ODME kalibrace. V každé iteraci se provede dopravní přiřazení, vyhodnotí se odchylky na kalibračních profilech a pomocí select-link proporcí se upraví O-D matice. Součástí této fáze jsou ochranné mechanismy, například omezení maximální změny poptávky, ořez vůči počáteční matici a návrat k nejlepšímu dosaženému stavu při zhoršování účelové funkce.

Poslední fáze slouží k jemnému doladění odchylek na úrovni screenlines. Používá nižší intenzitu korekcí než hlavní ODME smyčka, aby nedošlo ke zhoršení stability celkové matice. Výstupem kalibrace je upravená O-D matice a kalibrační report obsahující průběh metrik v jednotlivých iteracích.

Volitelným navazujícím krokem je \texttt{tune-supply}, který optimalizuje nabídkové parametry sítě (rychlostní a~kapacitní faktory po~třídách komunikací) pomocí koordinátového sestupu. Pro každou kombinaci faktorů se spustí zkrácená kalibrace a~vyhodnotí kompozitní účelová funkce zahrnující $R^2$, sklon regrese, \%RMSE a~screenline odchylky. Nejlepší nalezené faktory se aplikují na~síť a~následuje plná rekalibrace.

\subsection{Validační metriky}
\label{sec:impl-calibration-metrics}

Po každé iteraci kalibrace se vyhodnocuje sada metrik: GEH statistika, procentuální RMSE, koeficient determinace~$R^2$, sklon regresní přímky, MAPE, globální bias a~screenline odchylky. Protože GEH je primárně navržen pro hodinové objemy a~u~denních modelů je jeho vypovídací schopnost omezená, konvergenční rozhodnutí kombinuje více ukazatelů současně. Kalibrace se ukončuje po~splnění kompozitních kritérií, nebo po~stagnaci účelové funkce po~stanovený počet iterací (výchozí \textit{stall patience} je~8). Dosažené hodnoty jednotlivých metrik jsou vyhodnoceny v~kapitole~8.

Kalibrační a~validační data CSD jsou oddělena: kalibrace probíhá výhradně na~kalibrační podmnožině sčítacích profilů (poměr 65\,:\,35), zatímco krok \texttt{validate} používá nezávislou validační část.

% =========================================================================
\section{Implementace scénářů uzavírek}
\label{sec:impl-scenarios}

Klíčovou funkcionalitou systému je možnost vyhodnotit dopad dopravních uzavírek na modelovanou síť. Tato část popisuje implementační mechanismus od definice scénáře až po výpočet rozdílových metrik. Logika je realizována v modulech \texttt{scenarios/engine.py} a~\texttt{scenarios/state.py} a navazuje přímo na baseline model vytvořený v předchozích krocích pipeline.


\subsection{Definice scénáře}
\label{sec:impl-scenarios-definition}

Scénář je reprezentován jako množina modifikací aplikovaných na konkrétní hrany sítě identifikované pomocí \texttt{link\_id}. Každá položka definuje směr (\texttt{both}, \texttt{ab}, \texttt{ba}) a typ zásahu, který určuje způsob úpravy parametrů hrany.

Framework rozlišuje dva základní typy zásahů. Typ \texttt{full} reprezentuje úplnou uzavírku, kdy je kapacita nastavena na minimální hodnotu a cestovní čas výrazně penalizován, což vede k praktickému vyloučení hrany z trasování. Typ \texttt{lanes} modeluje redukci pruhů, kdy je kapacita škálována poměrem \texttt{lanes\_remaining} ku \texttt{lanes}, přičemž ostatní atributy zůstávají zachovány.

Scénář může být definován třemi způsoby: jako JSON payload přes REST API (\texttt{POST /api/scenarios/run}), jako YAML soubor ve složce \texttt{scenarios/}, nebo automaticky z aktuálních uzavírek pomocí modulu \texttt{scenarios/closures.py}. Vstupní data jsou validována pomocí Pydantic modelu, který kontroluje základní konzistenci (např. \texttt{link\_id}~$>0$, \texttt{lanes\_remaining}~$\leq$~\texttt{lanes}) a odstraňuje duplicity.

\begin{lstlisting}[language={},caption={Příklad JSON payloadu pro scénář uzavírky.},
  label={lst:scenario-json}]
{
  "links": [
    {
      "link_id": 12345,
      "direction": "both",
      "closure_type": "full",
      "lanes": 2,
      "lanes_remaining": 0
    },
    {
      "link_id": 12346,
      "direction": "ab",
      "closure_type": "lanes",
      "lanes": 3,
      "lanes_remaining": 1
    }
  ]
}
\end{lstlisting}


\subsection{Aplikace scénáře a přepočet}
\label{sec:impl-scenarios-lifecycle}

Zpracování scénáře je řízeno stavovým automatem implementovaným v modulu \texttt{scenarios/state.py}. Každý scénář je reprezentován objektem \texttt{ScenarioJob}, který prochází stavy \texttt{QUEUED}, \texttt{RUNNING} a koncovými stavy \texttt{SUCCEEDED}, \texttt{FAILED}, \texttt{CANCELLED} nebo \texttt{TIMED\_OUT}. Validita přechodů je vynucena funkcí \texttt{transition\_job}.

Po odeslání požadavku je scénář zařazen do fronty a zpracován asynchronně pomocí \texttt{ThreadPoolExecutor}. Samotný výpočet zahrnuje načtení baseline O-D matice, sestavení grafu z AequilibraE projektu a aplikaci modifikací pomocí funkce \texttt{apply\_scenario\_to\_graph}. Tato funkce upravuje kapacitu a cestovní čas přímo v paměťové reprezentaci grafu, aniž by docházelo ke změně topologie nebo identifikátorů hran.

Následně je spuštěno dopravní přiřazení se shodnou konfigurací jako v baseline výpočtu. Výsledky jsou serializovány do GeoJSON souboru ve složce \texttt{outputs/scenarios/} a zpřístupněny přes REST API.

Zásadní vlastností implementace je, že scénář nikdy nemodifikuje persistovaný baseline model. Veškeré úpravy probíhají pouze nad pracovní kopií grafu v paměti, která je po dokončení výpočtu zahozena. Tím je zajištěna reprodukovatelnost, izolace scénářů a přímá srovnatelnost výsledků.


\subsection{Rozdílová analýza}
\label{sec:impl-scenarios-delta}

Po dokončení přiřazení je provedena rozdílová analýza mezi baseline a scénářovým stavem. Funkce \texttt{\_build\_scenario\_geojson} spojuje výsledné objemy s geometrií sítě a pro každou hranu vypočítá základní delta ukazatele.

Pro každou hranu je určena absolutní změna objemu (\texttt{delta\_vol}), relativní změna (\texttt{delta\_pct}), změna poměru objem/kapacita (\texttt{delta\_voc}) a změna cestovního času (\texttt{delta\_ct}). Tyto hodnoty jsou uloženy spolu s původními baseline a scénářovými atributy v GeoJSON souboru.

Souhrnné metriky jsou poskytovány samostatným API endpointem, který vrací agregované informace, jako je počet významně ovlivněných hran, průměrná změna a extrémní hodnoty. Tyto výstupy jsou přímo využívány ve frontendové aplikaci, kde jsou rozdíly vizualizovány barevnou škálou (nárůst červeně, pokles zeleně).

\subsection{Validace scénáře}
\label{sec:impl-scenarios-validation}

Validace scénáře probíhá ve dvou vrstvách. První vrstva kontroluje samotný vstup požadavku pomocí Pydantic modelu \texttt{ScenarioRunRequest}. Ověřuje zejména strukturální konzistenci scénáře, například kladnou hodnotu \texttt{link\_id}, vztah \texttt{lanes\_remaining}~$\leq$~\texttt{lanes} a odstraňuje duplicitní dvojice \texttt{link\_id} a \texttt{direction}. Pokud některá z uvedených hran v aktuální síti neexistuje, daná modifikace je přeskočena a scénář může pokračovat s částečně platným vstupem.

Druhá vrstva zachycuje chyby vzniklé během samotného výpočtu. Pokud při aplikaci scénáře nebo přiřazení dojde k výjimce, například v důsledku nedostupného centroidu po rozpadu sítě, stavový automat převede úlohu do stavu \texttt{FAILED} a chybová zpráva je zpřístupněna uživateli přes status endpoint. Explicitní kontrola souvislosti sítě po aplikaci uzavírky v současné verzi implementována není. Jedná se o plánované rozšíření, které by umožnilo uživatele upozornit na nevalidní scénář ještě před spuštěním výpočetně náročného přiřazení.

% =========================================================================
\section{Webová aplikace a vizualizace výsledků}
\label{sec:impl-frontend}

Poslední část implementace tvoří webová aplikace, která zpřístupňuje výsledky modelu uživateli. Jejím účelem není provádět kompletní výpočet dopravního modelu od začátku, ale pracovat s artefakty vytvořenými v předchozích krocích pipeline a umožnit jejich interaktivní prohlížení. Serverová část je spuštěna krokem \texttt{serve} a poskytuje data prostřednictvím REST API. Klientská část je implementována jako mapová aplikace, která umožňuje zobrazit referenční stav, scénářové výsledky a rozdílové vrstvy dopadů uzavírek.

Architektura aplikace je rozdělena na backendovou a frontendovou část. Backend je implementován v jazyce Python pomocí frameworku FastAPI a zajišťuje načítání výstupních souborů, jejich transformaci do formátů vhodných pro webové rozhraní a obsluhu scénářových požadavků. Frontend je implementován pomocí knihovny React a jazyka TypeScript. Komunikace mezi oběma částmi probíhá přes JSON a GeoJSON endpointy, které jsou vhodné jak pro tabulková data, tak pro mapové vrstvy.

\subsection{REST API (FastAPI)}
\label{sec:impl-frontend-api}

REST API je implementováno v modulu \texttt{api.py}. Jeho úkolem je zpřístupnit webové aplikaci výstupy pipeline, reporty a scénářové funkce. Server je spouštěn pomocí Uvicornu na portu 8000 a obsahuje CORS middleware pro komunikaci s frontendovým vývojovým serverem.

API je rozděleno do několika tematických skupin podle typu poskytovaných dat. První skupinu tvoří síťová data, tedy hrany, uzly, zóny, centroidy, modelová oblast a základní metadata. Druhá skupina endpointů slouží pro práci se scénáři uzavírek, zejména pro spuštění scénáře, sledování jeho stavu a načtení výsledků. Další endpointy zpřístupňují aktivní uzavírky, reporty z kalibrace a validace, časové profily a diagnostické výstupy. Vybrané endpointy jsou uvedeny v tabulce~\ref{tab:api-endpoints}.

\begin{table}[ht]
\centering
\caption{Přehled vybraných REST API endpointů systému.}
\label{tab:api-endpoints}
\small
\begin{tabular}{llp{5.5cm}}
\toprule
\textbf{Metoda} & \textbf{Cesta} & \textbf{Popis} \\
\midrule
\multicolumn{3}{l}{\textit{Síťová data}} \\
GET & \texttt{/api/links} & GeoJSON hran s dopravními objemy a možností filtrování \\
GET & \texttt{/api/zones} & GeoJSON dopravních zón \\
GET & \texttt{/api/meta} & Základní metadata modelu a nastavení mapy \\
\midrule
\multicolumn{3}{l}{\textit{Scénáře}} \\
POST & \texttt{/api/scenarios/run} & Spuštění scénáře uzavírky \\
GET & \texttt{/api/scenarios/\{id\}/status} & Stav scénářového výpočtu \\
GET & \texttt{/api/scenarios/\{id\}/results} & Výsledný GeoJSON scénáře \\
GET & \texttt{/api/scenarios/\{id\}/delta-summary} & Souhrn rozdílové analýzy \\
\midrule
\multicolumn{3}{l}{\textit{Uzavírky a reporty}} \\
GET & \texttt{/api/closures} & Aktivní uzavírky pro zvolené datum \\
POST & \texttt{/api/closures/run-scenario} & Spuštění scénáře z existujících uzavírek \\
GET & \texttt{/api/reports/calibration} & Kalibrační report ve formátu JSON \\
GET & \texttt{/api/reports/validation} & Validační report \\
\midrule
\multicolumn{3}{l}{\textit{Diagnostika}} \\
GET & \texttt{/api/diagnostics/bias} & Prostorové zobrazení odchylky modelu \\
GET & \texttt{/api/diagnostics/corridors} & Koridorová analýza \\
GET & \texttt{/api/diagnostics/zone-route} & Trasa mezi vybranými zónami \\
\bottomrule
\end{tabular}
\end{table}

Důležitou částí backendu je načítání síťových dat. Hrany sítě jsou uloženy odděleně od výsledků přiřazení, a proto se při prvním požadavku spojuje geometrie ze souboru \texttt{network\_links.gpkg} s tabulkou \texttt{assignment\_results.parquet}. Výsledek je následně uložen do cache, aby se při dalších požadavcích nemusel opakovat nákladnější diskový přístup a slučování dat. Cache je invalidována v případě, že se na disku objeví novější soubor výsledků, například po opětovném spuštění přiřazení nebo kalibrace.

\subsection{Mapové rozhraní (React + Leaflet)}
\label{sec:impl-frontend-map}

Frontendová aplikace je implementována jako jednostránková aplikace (Single Page Application, SPA) pomocí knihovny React. Pro typovou kontrolu je použit TypeScript a sestavení aplikace zajišťuje nástroj Vite. Navigace je rozdělena do několika hlavních částí: mapové rozhraní, diagnostika, reporty a informační stránka o modelu.

Hlavní část aplikace tvoří mapa implementovaná pomocí knihovny React-Leaflet. Jednotlivé prvky mapy jsou rozděleny do samostatných komponent, což umožňuje oddělit logiku vykreslování silniční sítě, zón, centroidů, uzavírek a scénářových výsledků. Přehled hlavních mapových komponent uvádí tabulka~\ref{tab:map-components}.

\begin{table}[ht]
\centering
\caption{Vybrané React komponenty mapového rozhraní.}
\label{tab:map-components}
\begin{tabular}{lp{7cm}}
\toprule
\textbf{Komponenta} & \textbf{Funkce} \\
\midrule
\texttt{LinksLayer} & Zobrazení hran sítě a referenčních dopravních objemů \\
\texttt{ZonesLayer} & Zobrazení polygonů dopravních zón \\
\texttt{CentroidsLayer} & Zobrazení centroidů zón \\
\texttt{ClosuresLayer} & Zobrazení aktivních uzavírek pro vybrané datum \\
\texttt{DeltaLinksLayer} & Zobrazení rozdílů mezi referenčním a scénářovým stavem \\
\texttt{TimeSelector} & Výběr data a denní periody \\
\texttt{ViewToggle} & Přepínání mezi pohledy baseline, scenario a delta \\
\bottomrule
\end{tabular}
\end{table}

Stav aplikace je rozdělen do dvou logických částí. První část uchovává nastavení mapy, například aktivní vrstvy, zvolenou časovou periodu a parametry zobrazení. Druhá část uchovává stav scénářového výpočtu, tedy definici uzavírky, identifikátor spuštěné úlohy a výsledky po dokončení výpočtu. Pro načítání dat z API jsou použity hooky knihovny TanStack Query, které zjednodušují práci s asynchronními požadavky a umožňují automatické opakované dotazování na stav scénáře.

% \begin{figure}[ht]
% \centering
% \includegraphics[width=\textwidth]{img/frontend_screenshot.png}
% \caption{Webová aplikace -- mapové rozhraní s vizualizací dopravních intenzit na síti modelu Brna.}
% \label{fig:frontend-screenshot}
% \end{figure}

\subsection{Scénářový panel}
\label{sec:impl-frontend-scenario}

Práce se scénáři uzavírek je soustředěna do komponenty \texttt{ScenarioPanel.tsx}. Tato část rozhraní propojuje uživatelský výběr v mapě s backendovým scénářovým výpočtem. Uživatel nejprve vybere hrany, které mají být ve scénáři omezeny, a následně pro každou z nich nastaví typ zásahu, směr a případný počet zbývajících pruhů.

Po odeslání scénáře aplikace vytvoří požadavek na endpoint \texttt{/api/scenarios/run}. Backend vrátí identifikátor scénářové úlohy, podle kterého frontend periodicky dotazuje endpoint \texttt{/api/scenarios/\{id\}/status}. Po dokončení výpočtu jsou automaticky načteny výsledky scénáře a aplikace se přepne do rozdílového zobrazení. V tomto pohledu vrstva \texttt{DeltaLinksLayer} zvýrazňuje změny dopravních objemů vůči referenčnímu stavu.

Scénářový panel současně zobrazuje souhrn dopadů, který je získán z endpointu \texttt{/api/scenarios/\{id\}/delta-summary}. Uživatel tak kromě mapové vizualizace získá také číselný přehled, například počet významně ovlivněných hran, průměrnou absolutní změnu intenzity a seznam nejvíce dotčených úseků.

\subsection{Dashboard a diagnostika}
\label{sec:impl-frontend-dashboard}

Vedle hlavního mapového rozhraní obsahuje aplikace diagnostické a~reportovací pohledy. Diagnostická stránka nabízí prostorovou vizualizaci odchylky modelu od~pozorovaných dat, zobrazení podílu tranzitní dopravy, koridorovou analýzu vybraných silničních tahů a~porovnání tras mezi zónami. Reportovací stránka shrnuje kalibrační a~validační metriky formou interaktivních grafů a~popisuje strukturu O-D matice podle segmentů.

% Formáty výstupních dat jsou popsány v sekci~\ref{sec:impl-network-export}.

\section{Shrnutí implementace}
\label{sec:impl-summary}

Tato kapitola popsala implementaci dopravního simulačního frameworku od přípravy vstupních geodat až po webovou aplikaci pro vizualizaci výsledků. Implementované řešení je postaveno na modulární pipeline řízené skriptem \texttt{run.py}, která zahrnuje 21 navazujících kroků. Ty pokrývají import a normalizaci silniční sítě, tvorbu dopravních zón, generování poptávky, dopravní přiřazení, kalibraci, validaci a zpřístupnění výsledků přes REST API. Výsledný model města Brna obsahuje 52\,493 hran, 47\,181 uzlů a 219 dopravních zón, z toho 191 interních a 28 externích gateway zón. Celková odhadovaná denní poptávka činí 935\,706 vozidel.

Důležitou vlastností implementace je oddělení obecné metodiky od městsky specifických parametrů. Základní nastavení frameworku je definováno ve výchozích konfiguračních hodnotách, zatímco konkrétní model města je řízen souborem \texttt{sim.yaml} v adresáři \texttt{config/<city>/}. Tento přístup podporuje reprodukovatelnost výpočtu a zároveň usnadňuje přenositelnost řešení na další města. Stejná pipeline může být použita nad jinou modelovou oblastí bez zásahu do zdrojového kódu, pokud jsou doplněny odpovídající vstupní parametry a datové zdroje.

Implementovaná webová vrstva zpřístupňuje výsledky modelu prostřednictvím REST API a interaktivního mapového rozhraní. Backend umožňuje načítat síťová data, reporty, diagnostické výstupy a scénářové výsledky, zatímco frontend poskytuje uživateli mapové vrstvy referenčního stavu, scénářového stavu a rozdílové analýzy. Scénářové výpočty jsou realizovány tak, aby nemodifikovaly persistovaný baseline model. Úpravy uzavírek se aplikují pouze na pracovní kopii grafu v paměti, což zachovává srovnatelnost jednotlivých scénářů.

Součástí implementace je také kalibrační mechanismus zahrnující gravitační kalibraci, gateway korekce, ODME a screenline vyhodnocení. V této kapitole byl popsán především způsob jeho realizace a datové vazby na ostatní části pipeline. Konkrétní dosažené kalibrační a validační výsledky jsou dále analyzovány v kapitole~8.

Současná implementace má také několik omezení. Model je založen na statickém dopravním přiřazení, a proto nezachycuje dynamiku tvorby front ani časově proměnné chování kongescí. Kvalita výsledků je závislá na úplnosti a správnosti dat OpenStreetMap a na dostupnosti kalibračních dat. Poptávková část vychází zejména ze SLDB~2021, které představuje statický snímek dojížďkového chování a nezachycuje elastickou reakci poptávky na dopravní omezení. Tato omezení je nutné zohlednit při interpretaci scénářových výsledků.

Na implementovaný systém navazuje následující kapitola, která se zaměřuje na experimentální vyhodnocení. V ní jsou analyzovány kalibrační a validační metriky, citlivost modelu a dopady vybraných scénářů dopravních uzavírek.
%
% =============================================================================
% ── Kontrolní seznam pro revizi (z doporuceni.txt) ──
% =============================================================================
% [ ] Je v úvodu kap. 7 explicitně řečeno, že navazuje na kap. 6?
% [ ] Je každá část kap. 7 navázána na konkrétní část kap. 6
%     a neopakuje návrhovou teorii?
% [x] Odpovídá popis pipeline skutečnému STEPS z run.py (21 kroků)?
% [ ] Jsou tvrdé metriky/počty převzaty z doložitelných výstupů?
% [ ] Je jasně rozlišeno, co je framework default a co city override?
% [ ] Je popsán životní cyklus scénáře a delta analýza?
% [ ] Je samostatně popsána aplikační vrstva (REST API + frontend)?
% [ ] Přechází závěr kap. 7 přirozeně do kap. 8?
% [ ] Neuvádí text kalibrační výsledky, které patří do kap. 8?
% [ ] Je terminologie konzistentní (čeština + eng. v závorce)?
% =============================================================================
